% ============================================================================
\clearpage
\section{Decisión de posición y riesgo conjunto}\label{sec:decision}
% ============================================================================

\ficha{Pronósticos $\P$ calibrados por activo y horizonte, posición actual, costos, límites y estado de
confianza.}
{Construir escenarios conjuntos, optimizar rendimiento, CVaR y costos, y aplicar una zona de no
operación.}
{Región de pesos estables o acción discreta (mantener, aumentar, reducir), con su razón y su riesgo
conjunto.}
{Unidades y horizontes compatibles; ninguna recomendación sin pasar por límites y abstención.}

\subsection{Escenarios conjuntos físicos}\label{sec:escenarios}

El riesgo de una cartera depende de cómo se mueven juntos sus activos; la suma de riesgos individuales no
lo reproduce. Los escenarios conjuntos se construyen con un modelo de factores de mercado y sector,
$R_i=\alpha_i+\beta_i^\top f+\varepsilon_i$, y una covarianza regularizada por contracción hacia una matriz
estructurada \citep{ledoit2004}. Como competidor se usa dependencia $t$ \citep{demarta2005}, que produce
más movimientos extremos conjuntos con la misma correlación lineal
(figuras~\ref{fig:copula_gauss} y~\ref{fig:copula_t}), y escenarios de estrés en los que la dependencia es
distinta de la de calma (\figref{fig:calma_crisis}).

\grafica[0.54]{s6_copula_gauss}{Cópula gaussiana con correlación 0.5 (datos sintéticos). En rojo, los
escenarios en que ambos activos quedan por debajo de su cuantil 2.5\%.}{fig:copula_gauss}

\grafica[0.54]{s6_copula_t}{Cópula $t$ con tres grados de libertad y la misma correlación: la frecuencia de
caídas conjuntas extremas casi se duplica.}{fig:copula_t}

\grafica[0.54]{s6_calma_crisis}{Dos regímenes con dependencia distinta. Un modelo que impone la misma
dependencia en calma y en crisis subestima el riesgo conjunto justo cuando importa.}{fig:calma_crisis}

\subsection{El problema costo--riesgo--rendimiento}\label{sec:problema}

Para acciones y ETF, con pesos $w$, retorno aleatorio conjunto $R$ al horizonte de decisión y posición
actual $w_{\text{actual}}$, el problema inicial es
\begin{equation}
\min_{w\in\mathcal C_t}\ -\E_{\widehat P_t}\bigl[w^\top R\bigr]
+\lambda\,\CVaR_{\alpha,\widehat P_t}\bigl(-w^\top R\bigr)+c\bigl(w-w_{\text{actual}}\bigr).
\label{eq:objetivo}
\end{equation}
Con $S$ escenarios $r_s$ y la representación de \citet{rockafellar2000},
$\CVaR_\alpha(L)=\min_\zeta\ \zeta+\frac{1}{(1-\alpha)S}\sum_s(L_s-\zeta)^+$, y con costos proporcionales
$c(\Delta w)=\kappa^\top|\Delta w|$, el problema es un programa lineal:
\begin{equation}
\begin{aligned}
\min_{w,\,b,\,v,\,\zeta,\,z}\quad & -\hat\mu^\top w+\lambda\Bigl(\zeta+\frac{1}{(1-\alpha)S}\sum_{s=1}^{S}z_s\Bigr)+\kappa^\top(b+v)\\
\text{s.a.}\quad & z_s\ge -r_s^\top w-\zeta,\quad z_s\ge 0,\qquad s=1,\dots,S,\\
& w=w_{\text{actual}}+b-v,\quad b,v\ge 0,\\
& \mathbf 1^\top w+\kappa^\top(b+v)\le 1,\quad 0\le w\le w_{\max},\quad w\in\mathcal C_t .
\end{aligned}
\label{eq:lp}
\end{equation}
$\CVaR_{0.95}$ es el promedio del 5\% peor de las pérdidas simuladas (\figref{fig:perdidas}). El conjunto
factible $\mathcal C_t$ recoge las restricciones de la \tabref{tab:restricciones}. Rendimiento, riesgo y
costos deben expresarse en las mismas unidades y el mismo horizonte; los costos incluyen compra y venta,
fracciones de acción, mínimos por operación y, cuando corresponda, conversión de moneda.

\begin{table}[H]
\centering\small
\caption{Restricciones del conjunto factible $\mathcal C_t$.}
\label{tab:restricciones}
\begin{tabularx}{\linewidth}{@{}l Y@{}}
\toprule
\encab{Restricción} & \encab{Forma lineal típica}\\
\midrule
Presupuesto y caja & $\mathbf 1^\top w+\text{costos}\le 1$; sin apalancamiento en la primera versión\\
Concentración & $w_i\le w_{\max}$ por activo\\
Exposición sectorial & $\sum_{i\in\text{sector}}w_i\le s_{\max}$\\
Beta de la cartera & $\beta^\top w\le\beta_{\max}$\\
Rotación & $\mathbf 1^\top(b+v)\le\text{rot}_{\max}$ por sesión\\
Liquidez & valor negociado $\le$ fracción del volumen diario promedio\\
Presupuesto de riesgo (opcional) & $\CVaR_\alpha(-w^\top R)\le b_{\text{riesgo}}$\\
\bottomrule
\end{tabularx}
\end{table}

\begin{notatecnica}[Homogeneidad del objetivo]
Sin restricciones activas, el objetivo \eqref{eq:objetivo} con costos nulos es positivamente homogéneo:
duplicar $w$ duplica el rendimiento esperado y el CVaR. La solución tiende entonces a los extremos
---toda la caja o la frontera de $\mathcal C_t$--- y es muy sensible a $\lambda$. En un experimento
sintético de la implementación de referencia (prueba \texttt{test\_objetivo\_homogeneo\_lleva\_a\_extremos}),
con escenarios a cinco sesiones, $\lambda=0.01$ lleva a los límites de concentración y $\lambda=0.02$ a
caja total. Conviene fijar la escala con un presupuesto de riesgo
explícito (última fila de la \tabref{tab:restricciones}), que además es más fácil de interpretar.
\end{notatecnica}

\subsection{CVaR: qué mide y qué no}\label{sec:cvar}

\grafica[0.6]{s6_perdidas}{Distribución simulada de pérdidas a cinco sesiones de una cartera sintética de
cuatro activos. El VaR 95\% es el umbral del 5\% peor; el CVaR 95\% es el promedio de esa cola.}{fig:perdidas}

\grafica[0.56]{s6_suma_riesgos}{Suma ponderada de los CVaR individuales frente al CVaR de la cartera. La
diferencia, que es el beneficio de diversificación, casi desaparece en el escenario de estrés.}{fig:suma_riesgos}

\begin{noconcluir}
El CVaR evalúa el promedio de pérdidas en la cola especificada \emph{bajo los escenarios supuestos}. No
es un límite garantizado de pérdida: si los escenarios subestiman la dependencia o las colas, el CVaR
subestima el riesgo.
\end{noconcluir}

\subsection{Zona de no operación y estabilidad de la decisión}\label{sec:no_operacion}

Un cambio pequeño de peso debe justificar sus costos \emph{y} superar una banda de incertidumbre. El
rendimiento esperado es la entrada más ruidosa del problema; al remuestrear escenarios y perturbar
$\hat\mu$, el óptimo recorre buena parte de la frontera del presupuesto de riesgo
(\figref{fig:region_estable}). Por eso la salida no es un punto sino una \emph{región casi óptima}
\begin{equation}
\mathcal R_\delta=\bigl\{w\in\mathcal C_t:\ \hat\mu^\top w\ \ge\ \hat\mu^\top w^\ast-\delta\bigr\},
\label{eq:region}
\end{equation}
con $\delta$ del orden de la dispersión bootstrap del valor del óptimo. La regla es simple: si
$w_{\text{actual}}\in\mathcal R_\delta$, mantener; si no, moverse solo hasta la región, con el traslado de
menor costo $\arg\min_{w\in\mathcal R_\delta}\kappa^\top|w-w_{\text{actual}}|$, no hasta el punto
$w^\ast$. La recomendación se expresa como región o como acción discreta (mantener, aumentar, reducir,
cerrar), evitando falsa precisión como ``7.431\%'' cuando el rango razonable es amplio
(\figref{fig:histograma_pesos}).

\grafica[0.62]{s6_region_estable}{Región casi óptima en el plano de pesos de dos activos (supuesto: el
rendimiento esperado tiene una desviación estándar del 30\% de su valor). Los puntos son óptimos
remuestreados; la línea discontinua, el borde interior de la región. La cartera actual dentro de la región
se mantiene; la que excede el presupuesto se reduce solo hasta la región.}{fig:region_estable}

\grafica[0.58]{s6_histograma_pesos}{Distribución del peso óptimo en A entre remuestreos. La cifra puntual
aparenta una precisión que el rango del 10\% al 90\% desmiente.}{fig:histograma_pesos}

\subsection{Fase posterior: optimización robusta}\label{sec:dro}

En una fase posterior la distribución puntual $\widehat P_t$ se reemplaza por un conjunto de
distribuciones plausibles y se optimiza el peor caso dentro de él:
\begin{equation}
\min_{w\in\mathcal C_t}\ \sup_{Q\in\mathcal B_\varepsilon(\widehat P_t)}\ \E_Q\bigl[\ell(w,R)\bigr],
\qquad \mathcal B_\varepsilon(\widehat P_t)=\{Q:\ W(Q,\widehat P_t)\le\varepsilon\}.
\label{eq:dro}
\end{equation}
Una bola de Wasserstein es una opción natural (\figref{fig:ambiguedad}); para pérdidas lineales por
tramos, como la del CVaR, el problema admite reformulaciones tratables \citep{esfahani2018}. El radio se
selecciona con validación temporal y pruebas de sensibilidad, no a partir de un porcentaje de confianza
inventado.

\begin{figure}[H]
\centering
\resizebox{0.86\linewidth}{!}{% Conjunto de ambigüedad de Wasserstein alrededor de la distribución estimada.
\begin{tikzpicture}[x=1cm,y=1cm]
\fill[qaqua!8] (0,0) circle (2.1);
\draw[qaqua,line width=0.8pt] (0,0) circle (2.1);
\fill[tinta] (0,0) circle (2.2pt);
\node[font=\sffamily\scriptsize,anchor=north] at (0,-0.08) {$\widehat{P}_t$};
\foreach \a/\r in {20/1.2,75/0.8,130/1.5,200/1.1,250/1.7,310/0.9,345/1.6}
  \fill[tenue] (\a:\r) circle (1.6pt);
\fill[qnaranja] (35:2.1) circle (2.6pt);
\node[font=\sffamily\scriptsize,anchor=south west] at (35:2.15) {$Q^\ast(w)$: peor caso para la cartera $w$};
\draw[flecha,draw=tinta2] (0,0) -- (-145:2.1);
\node[nota,anchor=north east] at (-150:1.2) {radio $\varepsilon$};
\node[font=\sffamily\scriptsize,star,star points=5,fill=qazul,inner sep=1.4pt] at (-40:1.35) {};
\node[nota,anchor=west] at (-40:1.5) {$P$ verdadera (desconocida)};
\node[nota,anchor=west,text width=5.2cm,align=left] at (3.2,0.4) {\textbf{Conjunto de ambigüedad}\\[1pt]
  $\mathcal{B}_\varepsilon(\widehat P_t)=\{Q: W(Q,\widehat P_t)\le\varepsilon\}$.\\[3pt]
  $\varepsilon$ se elige por validación temporal y pruebas de sensibilidad: pequeño, optimiza para la
  estimación; grande, se vuelve excesivamente conservador.};
\end{tikzpicture}
}
\caption{Conjunto de ambigüedad de Wasserstein. La decisión se protege contra el peor caso dentro de la
bola, que puede no contener a la distribución verdadera.}
\label{fig:ambiguedad}
\end{figure}

\begin{noconcluir}
Las garantías teóricas de esta familia dependen de supuestos (por ejemplo, muestras independientes y
colas controladas) que no se trasladan automáticamente a mercados dependientes y cambiantes.
\end{noconcluir}

\subsection{Control predictivo con horizonte móvil}\label{sec:mpc}

Después se prueba el control predictivo: planear varias decisiones, ejecutar solo la primera y recalcular
con datos nuevos \citep{boyd2017} (\figref{fig:mpc}):
\begin{equation}
\max_{w_t,\dots,w_{t+H-1}}\ \sum_{k=0}^{H-1}\Bigl(\hat\mu_{t+k}^\top w_{t+k}-\lambda\,\mathcal{R}_{t+k}(w_{t+k})
-c\bigl(w_{t+k}-w_{t+k-1}\bigr)\Bigr).
\label{eq:mpc}
\end{equation}
Requiere escenarios temporales coherentes entre horizontes, no simplemente marginales de 1, 5 y 20 días
ensambladas.

\begin{figure}[H]
\centering
\resizebox{0.8\linewidth}{!}{% Control predictivo con horizonte móvil: planear varias decisiones, ejecutar solo la primera.
\begin{tikzpicture}[x=1.45cm,y=1cm,
  plan/.style={circle,draw=qaqua,fill=white,line width=0.8pt,minimum size=2.6mm,inner sep=0pt},
  ejec/.style={circle,draw=qaqua,fill=qaqua,line width=0.8pt,minimum size=3.2mm,inner sep=0pt}]
\node[titulo,anchor=east] at (-0.35,1.5) {Plan en $t$};
\node[titulo,anchor=east] at (-0.35,0) {Plan en $t+1$};
\foreach \k/\y in {0/1.75,1/1.45,2/1.3,3/1.25,4/1.2} \node[plan] (a\k) at (\k,\y) {};
\foreach \k in {0,...,3}{\pgfmathtruncatemacro{\n}{\k+1}\draw[qaqua!60,line width=0.6pt] (a\k) -- (a\n);}
\node[ejec] at (0,1.75) {};
\foreach \k/\y in {1/0.2,2/-0.05,3/-0.15,4/-0.2,5/-0.22} \node[plan] (b\k) at (\k,\y) {};
\foreach \k in {1,...,4}{\pgfmathtruncatemacro{\n}{\k+1}\draw[qaqua!60,line width=0.6pt] (b\k) -- (b\n);}
\node[ejec] at (1,0.2) {};
\draw[guia,line width=0.7pt] (-0.2,-0.75) -- (5.4,-0.75);
\foreach \k/\e in {0/t,1/t+1,2/t+2,3/t+3,4/t+4,5/t+5}{\draw[guia] (\k,-0.75) -- (\k,-0.83);
  \node[nota,anchor=north] at (\k,-0.85) {$\e$};}
\draw[flecha,draw=qnaranja] (0.35,1.05) .. controls (0.6,0.8) .. (0.85,0.5);
\node[nota,anchor=west,text=tinta] at (0.72,0.95) {llegan datos nuevos: se recalcula todo el plan};
\node[nota,anchor=west] at (5.55,1.2) {solo se ejecuta\\ el primer paso\\ (círculo lleno)};
\end{tikzpicture}
}
\caption{Horizonte móvil. En cada sesión se optimiza una trayectoria de pesos y solo se ejecuta su primer
paso.}
\label{fig:mpc}
\end{figure}

\subsection{Si se gestionan opciones}\label{sec:opciones}

Si la posición incluye opciones, el P\&L lineal se sustituye por la \emph{revaluación completa} bajo
escenarios conjuntos de subyacente, superficie, paso del tiempo y ejercicio o asignación. La aproximación
delta--gamma--vega sirve como diagnóstico, no como valoración fiable de toda opción en movimientos grandes
(figuras~\ref{fig:opcion_spot} y~\ref{fig:opcion_vol}).

\grafica[0.56]{s6_opcion_spot}{P\&L a un día de un call a 30 días en el dinero: revaluación completa frente
a aproximaciones por griegas. Para movimientos grandes, delta--gamma sobrestima la ganancia y, a la
baja, inventa una pérdida que el call no puede tener.}{fig:opcion_spot}

\grafica[0.56]{s6_opcion_vol}{El mismo call cuando la volatilidad sube al caer el precio. Añadir vega no
arregla la aproximación en los extremos.}{fig:opcion_vol}

\begin{pendiente}
Confirmar si la posición del usuario incluye o incluirá opciones antes de construir esta rama.
\end{pendiente}

\begin{criterio}
Riesgo factible bajo todas las restricciones y resultados estables en pruebas de estrés: la recomendación
no cambia de signo ante perturbaciones plausibles de escenarios, costos y $\hat\mu$.
\end{criterio}
